Machine-learning intrusion detectors for quantum key distribution (QKD) report near-perfect accuracy on simulated attacks, yet their benign data are mostly stationary and their attackers are not tuned to the monitored statistics. We ask what window-level telemetry of a three-intensity decoy-state BB84 link can reveal. A faked-state attacker who reproduces the honest photon-number yields and error probabilities leaves the joint law of all detection and error counts unchanged; every telemetry detector then operates at chance while the finite-key analysis outputs a key the attacker knows. For attacks inside the device model the decoy-state bound already protects the key. We propose a decoy-consistency statistic (DCS) that refits the drifting channel parameters in every window and tests only the relations that drift cannot change, and we evaluate it with five baselines on 24-hour streams under four drift scales with an exact finite-key audit of attacker knowledge. Over \nSeeds{} seeds, isolation forests and random forests calibrated on stationary data alarm in \fprIFstatDOne\% and \fprRFstatDOne\% of benign windows under nominal drift; domain-randomised training lowers this to \fprRFDRDOne\%. The DCS stays between \fprDCSDOne\% and \fprDCSDFour\% up to fourfold drift and flags naive and gain-only blinding in the first 10-second window. All detectors, the DCS included, score an AUC between \aucCalMin{} and \aucCalMax{} on the calibrated attack, and every key block produced under any blinding variant was fully known to the attacker. Telemetry monitoring can catch imperfectly calibrated attacks without drowning operators in false alarms, but it cannot replace hardware countermeasures against detector control.
